% ---------------------------------------------------------------------------
% hades.tex — estilo LaTeX para los PDF exportados desde Neovim
%
% Se inyecta con --include-in-header sobre la plantilla LaTeX por defecto de
% pandoc, así que funciona con cualquier versión de pandoc y con xelatex,
% lualatex o pdflatex.
%
% Color principal: #9B5DE5
% ---------------------------------------------------------------------------

\makeatletter

% ---------------------------------------------------------------------------
% Paleta
% ---------------------------------------------------------------------------
\usepackage{xcolor}
\definecolor{hadesprimary}{HTML}{9B5DE5}  % color principal
\definecolor{hadesdark}{HTML}{4B2A7A}     % títulos / texto destacado
\definecolor{hadesmid}{HTML}{7B3FC9}      % subtítulos
\definecolor{hadeslight}{HTML}{F6F1FD}    % fondos (código, citas)
\definecolor{hadesgray}{HTML}{5F5A6B}     % texto secundario
\definecolor{hadesnote}{HTML}{0096C7}
\definecolor{hadestip}{HTML}{00A884}
\definecolor{hadeswarning}{HTML}{E09F00}
\definecolor{hadesimportant}{HTML}{9B5DE5}
\definecolor{hadescaution}{HTML}{D0213B}

% ---------------------------------------------------------------------------
% Paquetes
% ---------------------------------------------------------------------------
\usepackage{etoolbox}
\usepackage{titlesec}
\usepackage{fancyhdr}
\usepackage[breakable,skins]{tcolorbox}

% ---------------------------------------------------------------------------
% Metadatos (título y subtítulo) para la portada y la cabecera
% ---------------------------------------------------------------------------
\gdef\hades@title{}
\gdef\hades@subtitle{}
\let\hades@origtitle\title
\renewcommand{\title}[1]{\hades@origtitle{#1}\gdef\hades@title{#1}}
% pandoc define \subtitle con \providecommand (después de este fichero) y
% quarto antes de él: en ambos casos acaba mandando esta definición
\providecommand{\subtitle}{}
\renewcommand{\subtitle}[1]{\gdef\hades@subtitle{#1}}

% ---------------------------------------------------------------------------
% Bloque de título
% ---------------------------------------------------------------------------
\renewcommand{\maketitle}{%
  \begingroup
  \thispagestyle{hadesplain}%
  \def\and{\unskip,\ }%
  \def\thanks##1{}%
  \noindent{\color{hadesprimary}\rule{\linewidth}{3pt}}\par
  \vspace{1.2em}
  \noindent{\fontsize{24}{30}\selectfont\bfseries\color{hadesdark}\hades@title\par}
  \ifx\hades@subtitle\@empty\else
    \vspace{0.6em}
    \noindent{\Large\color{hadesmid}\hades@subtitle\par}
  \fi
  \vspace{1.2em}
  \ifx\@author\@empty\else
    \noindent{\large\color{hadesgray}\@author\par}
  \fi
  \ifx\@date\@empty\else
    \vspace{0.2em}
    \noindent{\color{hadesgray}\@date\par}
  \fi
  \vspace{0.8em}
  \noindent{\color{hadesprimary}\rule{\linewidth}{0.8pt}}\par
  \vspace{1.8em}
  \endgroup
}

% ---------------------------------------------------------------------------
% Encabezados de sección
% ---------------------------------------------------------------------------
\titleformat{\section}
  {\Large\bfseries\color{hadesdark}}
  {\color{hadesprimary}\thesection}{0.8em}{}
  [{\vspace{-0.6ex}\color{hadesprimary}\titlerule[0.8pt]}]
\titleformat{\subsection}
  {\large\bfseries\color{hadesmid}}
  {\thesubsection}{0.8em}{}
\titleformat{\subsubsection}
  {\normalsize\bfseries\color{hadesdark}}
  {\thesubsubsection}{0.8em}{}
\titlespacing*{\section}{0pt}{2.2ex plus .5ex minus .2ex}{1.4ex plus .2ex}
\titlespacing*{\subsection}{0pt}{1.8ex plus .5ex minus .2ex}{1ex plus .2ex}

% ---------------------------------------------------------------------------
% Cabecera y pie de página
% ---------------------------------------------------------------------------
\fancypagestyle{hadesplain}{%
  \fancyhf{}%
  \renewcommand{\headrulewidth}{0pt}%
  \renewcommand{\footrulewidth}{0pt}%
  \fancyfoot[C]{\small\color{hadesgray}\thepage}%
}
\fancypagestyle{hades}{%
  \fancyhf{}%
  \fancyhead[L]{\small\color{hadesgray}\hades@title}%
  \fancyhead[R]{\small\color{hadesgray}\nouppercase{\leftmark}}%
  \fancyfoot[C]{\small\color{hadesgray}\thepage}%
  \renewcommand{\headrulewidth}{0.4pt}%
  \renewcommand{\footrulewidth}{0pt}%
  \renewcommand{\headrule}{{\color{hadesprimary}%
    \hrule width\headwidth height\headrulewidth \vskip-\headrulewidth}}%
}
\setlength{\headheight}{14pt}
\pagestyle{hades}
% \tableofcontents usa el estilo plain: lo cambiamos por el nuestro
\fancypagestyle{plain}{%
  \fancyhf{}%
  \renewcommand{\headrulewidth}{0pt}%
  \fancyfoot[C]{\small\color{hadesgray}\thepage}%
}

% ---------------------------------------------------------------------------
% Listas
% ---------------------------------------------------------------------------
\renewcommand{\labelitemi}{\textcolor{hadesprimary}{\textbullet}}
\renewcommand{\labelitemii}{\textcolor{hadesmid}{\normalfont\bfseries\textendash}}
\renewcommand{\labelitemiii}{\textcolor{hadesmid}{\textasteriskcentered}}

% ---------------------------------------------------------------------------
% Bloques de código (solo existe Shaded si el documento tiene código resaltado)
% ---------------------------------------------------------------------------
\@ifundefined{Shaded}{}{%
  \renewenvironment{Shaded}{%
    \begin{tcolorbox}[breakable, enhanced, sharp corners, frame hidden,
      colback=hadeslight, boxrule=0pt,
      borderline west={2.5pt}{0pt}{hadesprimary},
      left=8pt, right=6pt, top=4pt, bottom=4pt,
      before skip=8pt, after skip=8pt]%
  }{%
    \end{tcolorbox}%
  }%
}

% ---------------------------------------------------------------------------
% Citas (> ...)
% ---------------------------------------------------------------------------
\renewenvironment{quote}{%
  \begin{tcolorbox}[breakable, enhanced, sharp corners, frame hidden,
    colback=hadeslight, boxrule=0pt,
    borderline west={3pt}{0pt}{hadesprimary},
    left=10pt, right=8pt, top=6pt, bottom=6pt,
    before skip=8pt, after skip=8pt, fontupper=\itshape\color{hadesgray}]%
}{%
  \end{tcolorbox}%
}

% ---------------------------------------------------------------------------
% Callouts (los genera callouts.lua):
%   \begin{hadescallout}{color}{Título} ... \end{hadescallout}
% ---------------------------------------------------------------------------
\newtcolorbox{hadescallout}[2]{%
  breakable, enhanced, sharp corners,
  colback=#1!6!white, colframe=#1,
  boxrule=0pt, leftrule=3pt, toprule=0pt, bottomrule=0pt, rightrule=0pt,
  colbacktitle=#1!14!white, coltitle=#1!70!black,
  title={\bfseries #2},
  left=8pt, right=8pt, top=4pt, bottom=4pt,
  before skip=10pt, after skip=10pt
}

% Callouts nativos de Quarto: mismos colores
\AtBeginDocument{%
  \definecolor{quarto-callout-note-color}{HTML}{0096C7}%
  \definecolor{quarto-callout-note-color-frame}{HTML}{0096C7}%
  \definecolor{quarto-callout-tip-color}{HTML}{00A884}%
  \definecolor{quarto-callout-tip-color-frame}{HTML}{00A884}%
  \definecolor{quarto-callout-important-color}{HTML}{9B5DE5}%
  \definecolor{quarto-callout-important-color-frame}{HTML}{9B5DE5}%
  \definecolor{quarto-callout-warning-color}{HTML}{E09F00}%
  \definecolor{quarto-callout-warning-color-frame}{HTML}{E09F00}%
  \definecolor{quarto-callout-caution-color}{HTML}{D0213B}%
  \definecolor{quarto-callout-caution-color-frame}{HTML}{D0213B}%
}

% ---------------------------------------------------------------------------
% Tablas
% ---------------------------------------------------------------------------
\@ifpackageloaded{colortbl}{}{\usepackage{colortbl}}
\arrayrulecolor{hadesprimary}

% ---------------------------------------------------------------------------
% Rótulos en castellano (si se define `lang`, babel los sobrescribe)
% ---------------------------------------------------------------------------
\renewcommand{\figurename}{Figura}
\renewcommand{\tablename}{Tabla}

\makeatother
